% docs/slides/preamble.tex -- 五份 deck 共享
\documentclass[aspectratio=169,11pt]{beamer}

% ==== 中文与字体 ====
\usepackage[UTF8,fontset=none]{ctex}
\usepackage{fontspec}

\IfFontExistsTF{PingFang SC}{%
  \setCJKmainfont{PingFang SC}[BoldFont=PingFang SC,ItalicFont=PingFang SC,AutoFakeSlant=0.15]
  \setCJKsansfont{PingFang SC}[ItalicFont=PingFang SC,AutoFakeSlant=0.15]
  \setCJKmonofont{PingFang SC}
}{%
  \IfFontExistsTF{Source Han Sans SC}{%
    \setCJKmainfont{Source Han Sans SC}
    \setCJKsansfont{Source Han Sans SC}
  }{%
    \IfFontExistsTF{Noto Sans CJK SC}{%
      \setCJKmainfont{Noto Sans CJK SC}
      \setCJKsansfont{Noto Sans CJK SC}
    }{%
      \setCJKmainfont{Heiti SC}
      \setCJKsansfont{Heiti SC}
    }
  }
}

\IfFontExistsTF{Helvetica Neue}{\setmainfont{Helvetica Neue}}{%
  \IfFontExistsTF{Inter}{\setmainfont{Inter}}{}}
\IfFontExistsTF{JetBrains Mono}{\setmonofont{JetBrains Mono}[Scale=0.85]}%
  {\setmonofont{Menlo}[Scale=0.85]}

% ==== 颜色 ====
\usepackage{xcolor}
\definecolor{cMain}{HTML}{1F4E79}
\definecolor{cAccent}{HTML}{E07A1F}
\definecolor{cGray}{HTML}{595959}
\definecolor{cMute}{HTML}{F4F4F4}
\definecolor{cOK}{HTML}{2E7D32}
\definecolor{cBad}{HTML}{B23A3A}
\definecolor{cWarn}{HTML}{C8A300}

% ==== 主题 ====
\usepackage{tcolorbox}
\usetheme{metropolis}
\metroset{progressbar=frametitle,numbering=fraction,block=fill}
\setbeamercolor{frametitle}{bg=cMain,fg=white}
\setbeamercolor{progress bar}{fg=cAccent,bg=cMute}
\setbeamercolor{alerted text}{fg=cAccent}
\setbeamercolor{title}{fg=cMain}

% metropolis 默认尝试 Fira Sans，未安装时 fallback 到 Latin Modern Sans。
% 这里在主题加载后显式覆盖，强制使用 macOS 上一定存在的 Helvetica Neue，
% 跟中文 PingFang 风格一致；找不到则按顺序尝试 Inter / Arial。
\IfFontExistsTF{Helvetica Neue}{%
  \setsansfont{Helvetica Neue}%
}{%
  \IfFontExistsTF{Inter}{\setsansfont{Inter}}{%
    \IfFontExistsTF{Arial}{\setsansfont{Arial}}{}}%
}

% ==== TikZ ====
\usepackage{tikz}
\usetikzlibrary{positioning,arrows.meta,fit,shapes.geometric,calc,%
  decorations.pathmorphing,backgrounds}

\tikzset{
  node-box/.style={draw=cMain,thick,rounded corners=2pt,
    minimum height=8mm,minimum width=18mm,align=center,font=\small},
  node-state/.style={draw=cMain,thick,circle,minimum size=14mm,
    align=center,font=\small},
  node-ok/.style={draw=cOK,thick,fill=cOK!10},
  node-warn/.style={draw=cWarn,thick,fill=cWarn!15},
  node-bad/.style={draw=cBad,thick,fill=cBad!10},
  % 色盲友好的状态机三态：蓝 (正常/Closed) / 橙 (告警/Open) / 灰 (中间/HalfOpen)
  % 避免 red-green 二元区分，改用 hue + lightness 双通道。
  node-state-normal/.style={draw=cMain,thick,fill=cMain!12},
  node-state-alert/.style={draw=cAccent,thick,fill=cAccent!18},
  node-state-transition/.style={draw=cGray,thick,fill=cGray!18},
  arrow/.style={-{Stealth[length=2mm]},thick,draw=cMain},
  arrow-async/.style={-{Stealth[length=2mm]},thick,dashed,draw=cGray},
  arrow-bad/.style={-{Stealth[length=2mm]},thick,draw=cBad},
  arrow-alert/.style={-{Stealth[length=2mm]},thick,draw=cAccent},
  lifeline/.style={dashed,thick,draw=cGray},
}

% ==== 代码 ====
\usepackage{listings}
\lstdefinelanguage{Go}{
  morekeywords={break,case,chan,const,continue,default,defer,else,
    fallthrough,for,func,go,goto,if,import,interface,map,package,range,
    return,select,struct,switch,type,var,nil,true,false,iota},
  morekeywords=[2]{string,int,int64,uint,uint64,bool,byte,error,context},
  sensitive=true,
  morecomment=[l]//,morecomment=[s]{/*}{*/},
  morestring=[b]",morestring=[b]`,
}
\lstdefinelanguage{LuaRedis}{
  morekeywords={and,break,do,else,elseif,end,false,for,function,if,in,
    local,nil,not,or,repeat,return,then,true,until,while},
  morekeywords=[2]{redis,call,KEYS,ARGV,tonumber,HGET,HSET,GET,SET,
    DECR,DECRBY,INCR,INCRBY,EXPIRE,PEXPIRE,ZADD,ZCARD,ZREMRANGEBYSCORE,
    EVAL,EVALSHA},
  sensitive=true,
  morecomment=[l]{--},
  morestring=[b]",morestring=[b]',
}

\lstdefinestyle{base}{
  basicstyle=\ttfamily\scriptsize,
  numbers=left,numberstyle=\tiny\color{cGray},
  keywordstyle=\color{cMain}\bfseries,
  keywordstyle=[2]\color{cAccent},
  stringstyle=\color{cOK},
  commentstyle=\color{cGray}\itshape,
  backgroundcolor=\color{cMute},
  frame=single,framerule=0pt,
  showstringspaces=false,
  breaklines=true,
  tabsize=2,
  columns=fullflexible,
}
\lstset{style=base}

% ==== 三线表与附录页码 ====
\usepackage{booktabs}
\usepackage{appendixnumberbeamer}

% ==== 页脚来源标注 ====
\newcommand{\srcnote}[1]{{\color{cGray}\scriptsize\itshape #1}}

% 关键论点框
\newcommand{\keypoint}[1]{%
  \begin{tcolorbox}[colback=cMute,colframe=cMain,boxrule=0.5pt,arc=2pt]%
  #1\end{tcolorbox}}


\title{AI 检索的业务视角}
\subtitle{ES 关键词 + Milvus 语义向量的 hybrid 路线}
\author{gomall · 业务侧 deck}
\date{}

\begin{document}

\maketitle

\begin{frame}{今日大纲}
\tableofcontents
\end{frame}

% ============================================================
\section{deck 12 的遗留问题}
% ============================================================

\begin{frame}{deck 12 给搜索做了什么}
\begin{itemize}
  \item LIKE $\to$ ES：解决了"分词 / 字段加权 / 相关性"三件事。
  \item Outbox + indexer：解决了"上架 5 分钟内可搜"的承诺。
  \item ES 挂 $\to$ DB 兜底：解决了"搜索是流量入口不能 5xx"的承诺。
\end{itemize}
\vspace{2mm}
\keypoint{deck 12 把"搜得到"做到位了。这份 deck 要回答的是：\textbf{搜不到的那部分} query 怎么办。}
\srcnote{repository/es/product\_index.go:156--209 当前 multi\_match 查询}
\end{frame}

\begin{frame}{ES 关键词解决不了的三类 query}
\begin{tabular}{@{}lll@{}}
\toprule
query 类型 & 例子 & ES 行为 \\
\midrule
自然语言   & \textit{送女朋友的礼物}    & 拆成"送 / 女朋友 / 礼物" $\to$ 命中乱 \\
跨表达同义 & \textit{苹果手机最新款}    & "苹果"+"手机"+"最新款" $\to$ 漏 iPhone 15 \\
口语描述   & \textit{适合健身的运动鞋} & 长尾词命中率极低 $\to$ 趋近 0 命中 \\
\bottomrule
\end{tabular}
\vspace{2mm}
\begin{itemize}
  \item ES 衡量的是\textbf{词的重合度}，不是\textbf{意思的接近度}。
  \item 用户已经习惯用自然语言搜（短视频 / 大模型养出来的）。
  \item 关键词召回率掉到 60\% 以下，长尾 GMV 直接漏走。
\end{itemize}
\end{frame}

\begin{frame}{为什么长尾搜索是大类目电商的命门}
\begin{itemize}
  \item 头部 SKU 的搜索词高度集中：\textit{iPhone}、\textit{口红}、\textit{洗衣液}。
  \item 但全站 GMV 里，头部 SKU 通常只贡献 30\%。
  \item 剩下 70\% 来自\textbf{长尾 SKU}：小众品类、个性化需求、口语化搜索。
  \item 业界口径：亚马逊有约 70\% GMV 来自非头部商品。
\end{itemize}
\vspace{2mm}
\keypoint{长尾搜索 miss 一次 = 一个细分类目的下单机会被默默放掉，
没有 5xx，没有告警，但 GMV 一直在漏。}
\end{frame}

% ============================================================
\section{三层检索体系}
% ============================================================

\begin{frame}{从单层 ES 到三层检索的整体形态}
\begin{tikzpicture}[node distance=4mm and 7mm,every node/.append style={font=\tiny}]
  \node[node-box]                       (q)   {用户\\query};
  \node[node-box,right=of q,node-warn]  (kw)  {关键词层\\ES};
  \node[node-box,right=of kw,node-warn] (sem) {语义层\\Milvus};
  \node[node-box,right=of sem]          (bz)  {业务调权\\运营 boost};
  \node[node-box,right=of bz,node-ok]   (mix) {融合排序\\RRF / 加权};
  \node[node-box,right=of mix]          (top) {top\_k};
  \draw[arrow] (q)   -- (kw);
  \draw[arrow] (kw)  -- (sem);
  \draw[arrow] (sem) -- (bz);
  \draw[arrow] (bz)  -- (mix);
  \draw[arrow] (mix) -- (top);
\end{tikzpicture}
\vspace{2mm}
\begin{itemize}
  \item 三层不是叠加，是\textbf{并行召回} + \textbf{统一融合}。
  \item 关键词召回保证"精确匹配不掉"，语义召回保证"长尾不漏"。
  \item 业务调权层让运营在不动算法的前提下调结果。
\end{itemize}
\srcnote{service/search/service.go:11--16 当前只走 ES 一路；融合在 feat/semantic-search PR}
\end{frame}

\begin{frame}{三层的业务责任划分}
\begin{tabular}{@{}lp{4.0cm}p{5.0cm}@{}}
\toprule
层 & 解决什么 & gomall 现状 \\
\midrule
关键词层 & \textit{iPhone 13} $\to$ name 命中 & 已在 main \\
语义层   & \textit{送礼物} $\to$ 礼物 / 饰品   & feat/milvus-vector + feat/semantic-search PR \\
业务调权 & 大促主推强制置顶 & 路线图，运营需求成熟后做 \\
\bottomrule
\end{tabular}
\vspace{2mm}
\begin{itemize}
  \item 关键词层："搜得到品牌型号"的下限。
  \item 语义层："搜得到长尾意图"的上限。
  \item 业务调权层："运营能干预"的政策接口。
\end{itemize}
\srcnote{repository/es/product\_index.go:17 product 索引；语义层 collection 在 feat/milvus-vector}
\end{frame}

% ============================================================
\section{hybrid 加权的业务含义}
% ============================================================

\begin{frame}{两路分数为什么不能直接加}
\begin{itemize}
  \item ES \texttt{\_score}：BM25 + 字段加权，绝对值几十到几百，跟 query 长度强相关。
  \item Milvus 向量距离：cosine 在 $[-1, 1]$，L2 在 $[0, +\infty)$，跟 query 长度无关。
  \item 直接加 $\to$ ES 分数把向量分数压死 $\to$ 等于没引入语义层。
\end{itemize}
\vspace{2mm}
\keypoint{不同量纲必须先 normalize。这是一个工程必修，但业务侧也要懂：
否则"上线了向量搜索召回率没变化"就解释不清楚。}
\srcnote{融合公式与 minMaxNormalize 实现在 feat/semantic-search PR 的 service/search/semantic.go}
\end{frame}

\begin{frame}{50/50 默认权重的业务判断}
\begin{tabular}{@{}lll@{}}
\toprule
权重 (kw : sem) & 适用场景 & 业务侧理由 \\
\midrule
70 : 30 & 商品库已全名上架  & 关键词命中率本来就高 \\
50 : 50 & 出厂默认           & 两路同等信任，先保召回 \\
30 : 70 & 自然语言 query 多 & 用户更多用描述 / 长尾词 \\
\bottomrule
\end{tabular}
\vspace{2mm}
\begin{itemize}
  \item 50/50 不是"算法最优"，是\textbf{产品对未知用户的初始信任}。
  \item 真正调权要靠埋点：人工标注或点击日志做 NDCG / MRR 离线评估。
  \item 也可以做 AB：按 UV 分桶切流量，观察 7 日 GMV / 0 结果率。
\end{itemize}
\end{frame}

\begin{frame}{融合方法：加权 vs RRF}
\begin{tikzpicture}[node distance=5mm and 10mm,every node/.append style={font=\tiny}]
  \node[node-box]                  (esh) {ES 100 条\\BM25 分数};
  \node[node-box,below=of esh]     (vch) {Milvus 100 条\\cosine 分数};
  \node[node-box,right=of esh,node-warn] (norm1) {min-max\\归一};
  \node[node-box,right=of vch,node-warn] (norm2) {min-max\\归一};
  \node[node-box,right=2cm of norm1,yshift=-7mm,node-ok] (mix) {加权融合\\$\alpha \cdot kw + (1-\alpha) \cdot sem$};
  \node[node-box,right=of mix]     (top) {top 20};
  \draw[arrow] (esh) -- (norm1);
  \draw[arrow] (vch) -- (norm2);
  \draw[arrow] (norm1) -- (mix);
  \draw[arrow] (norm2) -- (mix);
  \draw[arrow] (mix) -- (top);
\end{tikzpicture}
\vspace{2mm}
\begin{itemize}
  \item gomall 当前路线：加权 + min-max。
  \item RRF（reciprocal rank fusion）是另一条路：只看排名 \texttt{1/(rank+60)}，不看分数。
  \item 加权更可调，RRF 更稳健。前期选加权，留着 AB 调 $\alpha$。
\end{itemize}
\srcnote{feat/semantic-search PR 中 service/search/semantic.go 的 50/50 融合公式}
\end{frame}

% ============================================================
\section{语义层落地：embedding + 向量库}
% ============================================================

\begin{frame}{embedding 不只是"调一个 API"}
\begin{tabular}{@{}llll@{}}
\toprule
模型 & 维度 & 中文 & 成本 \\
\midrule
OpenAI text-embedding-3-small & 1536 & 一般  & \$0.02 / 1M tokens \\
OpenAI text-embedding-3-large & 3072 & 中等  & \$0.13 / 1M tokens \\
BGE-base-zh-v1.5               & 768  & 强   & 自部署 \\
BGE-large-zh-v1.5              & 1024 & 最强 & 自部署，GPU \\
GTE / E5-large-zh              & 1024 & 强   & 自部署 \\
\bottomrule
\end{tabular}
\vspace{2mm}
\begin{itemize}
  \item 中文电商：BGE-zh 系列在中文 benchmark 上长期领先 OpenAI。
  \item 自部署成本 = GPU 显存 + 推理服务 SRE。
  \item 调外部 API = 数据出境 + 单调用 50--500ms 延迟。
\end{itemize}
\end{frame}

\begin{frame}{gomall 的 embedding 抽象层}
\begin{itemize}
  \item 不绑死任何具体模型。env 切：
  \begin{itemize}
    \item \texttt{EMBEDDING\_API\_URL} 未设 $\to$ SHA-256 衍生 768 维占位向量，本地可跑。
    \item 设了 $\to$ POST \texttt{\{model, input\}}，OpenAI / BGE / Ollama 三种响应结构都兼容。
    \item \texttt{EMBEDDING\_MODEL} 切模型，\texttt{EMBEDDING\_API\_KEY} 走 Bearer。
  \end{itemize}
  \item 切模型只改 env，不动代码 —— 12-factor。
  \item 超时 5s，避免拖死检索 p99。
\end{itemize}
\vspace{2mm}
\srcnote{feat/semantic-search PR 的 service/search/embedding.go EmbedText 抽象}
\end{frame}

\begin{frame}{为什么不是"在 ES 里加个 dense\_vector 就行"}
\begin{itemize}
  \item ES 7.x 后支持 \texttt{dense\_vector}，但没有专门的 ANN 索引（早期版本）。
  \item 商品库 100 万行 + 768 维向量 = $\sim$3 GB 原始数据。
  \item ES 直接搜向量是\textbf{暴力近邻}，p99 退化到秒级。
  \item Milvus 用 HNSW 索引，同样 100 万向量 p99 在 5--20ms。
\end{itemize}
\vspace{2mm}
\keypoint{专用向量库不是"高大上"，是\textbf{ms 级响应的工程下限}。
搜索接口 p99 不能超过 200ms 是产品红线。}
\srcnote{repository/es/product\_index.go:17--40 ES product 索引故意不放向量字段}
\end{frame}

\begin{frame}{HNSW 还是 IVF：业务给出的答案}
\begin{tabular}{@{}llll@{}}
\toprule
索引 & 准确率 & 内存 & 适用规模 \\
\midrule
HNSW (M=16) & 高 ($\geq$95\%) & 高 ($\sim$2$\times$ 原始) & $<$ 1000 万 \\
IVF\_FLAT   & 中             & 低                       & 千万级 \\
IVF\_PQ     & 中低           & 极低                     & 亿级 \\
\bottomrule
\end{tabular}
\vspace{3mm}
\begin{itemize}
  \item gomall 商品数 $\sim$ 100 万，HNSW 直接选定。
  \item M=16 / efConstruction=200 / efSearch=64 是 Milvus 官方推荐的中等配置。
  \item 召回不达标 $\to$ 调 efSearch（运行时参数，不重建索引）。
\end{itemize}
\srcnote{feat/milvus-vector PR 的 repository/milvus/product\_vector.go HNSW 参数}
\end{frame}

\begin{frame}{静默降级：Milvus 挂了用户感受是什么}
\begin{tikzpicture}[node distance=8mm and 12mm]
  \node[node-state,node-state-normal] (n) {正常\\三层全开};
  \node[node-state,node-state-transition,right=of n] (d) {语义降级\\走 ES 单路};
  \node[node-state,node-state-alert,right=of d] (a) {全降级\\DB LIKE 兜底};
  \draw[arrow]       (n) -- node[above,font=\tiny]{Milvus 超时} (d);
  \draw[arrow-alert] (d) -- node[above,font=\tiny]{ES 也挂} (a);
  \draw[arrow,bend left=30] (d.south) to node[below,font=\tiny]{Milvus 恢复} (n.south);
\end{tikzpicture}
\vspace{2mm}
\begin{itemize}
  \item Milvus 客户端启动失败 $\to$ \texttt{tryInitMilvus} recover $\to$ 进程照常起。
  \item \texttt{MILVUS\_ADDR} 未设 $\to$ 直接跳过初始化（本地开发场景）。
  \item 语义召回挂 $\to$ ES 关键词单路仍可用 $\to$ 长尾词 miss 但精确词不掉。
\end{itemize}
\srcnote{cmd/main.go:60--69 tryInitES 的同套路；tryInitMilvus 在 feat/milvus-vector PR}
\end{frame}

\begin{frame}{增量更新：复用 outbox 不再造一遍轮子}
\begin{tikzpicture}[node distance=3mm and 3mm,every node/.append style={minimum width=12mm,font=\tiny}]
  \node[node-box]                       (api)  {商家\\改名};
  \node[node-box,right=of api]          (db)   {MySQL\\product};
  \node[node-box,right=of db,node-warn] (ob)   {outbox};
  \node[node-box,right=of ob]           (mq)   {RMQ\\product.\\changed};
  \node[node-box,right=of mq,node-ok]   (esi)  {ES\\indexer};
  \node[node-box,below=of esi,node-ok]  (mvi)  {Milvus\\indexer};
  \node[node-box,right=of esi]          (es)   {ES doc};
  \node[node-box,right=of mvi]          (mv)   {Milvus\\vector};
  \draw[arrow] (api) -- (db);
  \draw[arrow] (db)  -- (ob);
  \draw[arrow-async] (ob) -- (mq);
  \draw[arrow] (mq)  -- (esi);
  \draw[arrow] (mq)  |- (mvi);
  \draw[arrow] (esi) -- (es);
  \draw[arrow] (mvi) -- (mv);
\end{tikzpicture}
\vspace{2mm}
\begin{itemize}
  \item ES indexer 和 Milvus indexer 订阅\textbf{同一个事件}：\texttt{product.changed}。
  \item Milvus indexer 多一步：调 embedding 服务 $\to$ 写 vector。
  \item 整体 SLA 跟 ES 一致：上架 5 分钟内可搜（含语义召回）。
\end{itemize}
\srcnote{service/search/indexer.go:17--49 ES indexer 模板；Milvus indexer 在 feat/milvus-vector PR}
\end{frame}

% ============================================================
\section{业务收益与运营接口}
% ============================================================

\begin{frame}{四类 query 的命中对比（模拟）}
\begin{tabular}{@{}lllr@{}}
\toprule
query 类型 & 例子                          & kw only & hybrid 50/50 \\
\midrule
精确 SKU    & \textit{iPhone 13 Pro 256GB} & top1 准 & top1 准 \\
同义跨表达 & \textit{苹果手机最新款}        & miss    & top3 命中 \\
自然语言   & \textit{适合健身的运动鞋}     & 全 miss & top5 命中 \\
模糊场景   & \textit{送男朋友的礼物}       & 凌乱    & 相关性高 \\
\bottomrule
\end{tabular}
\vspace{3mm}
\begin{itemize}
  \item 精确 SKU：语义层无收益，但也不会拖低。
  \item 同义 / 长尾：从 0 命中提到 top3 $\to$ 这就是\textbf{增量 GMV} 的来源。
  \item 数字仅作业务侧示意，实际值要按上线后埋点测算。
\end{itemize}
\end{frame}

\begin{frame}{压测视角下的检索接口预算}
\begin{tabular}{@{}lrr@{}}
\toprule
链路 & RPS & p95 \\
\midrule
\texttt{/ping} 基线           & 64{,}254 & 3.51ms \\
\texttt{/product/show} (DB)   & 62{,}226 & 3.01ms \\
\texttt{/product/search} (ES) & --       & 业务承诺 $<$ 200ms \\
\texttt{/product/semantic} (ES + Milvus) & -- & 业务承诺 $<$ 300ms \\
\bottomrule
\end{tabular}
\vspace{2mm}
\begin{itemize}
  \item 详情页 62K RPS / p95 3ms 是基线，搜索接口本来就慢一个数量级。
  \item 引入语义层后 p99 预算放宽到 300ms：embedding 50--100ms + Milvus 20ms + ES 30ms + 融合 + DB 取详情。
  \item 业务承诺写在 FAQ：\textbf{搜索接口慢于 1s 视为故障}。
\end{itemize}
\srcnote{stressTest/REPORT.md 链路压测表 基线 64{,}254 RPS / product\_show 62{,}226 RPS}
\end{frame}

\begin{frame}{运营 boost：算法做不到的事 + 路线图}
\begin{itemize}
  \item \textbf{大促主推}：双 11 期间，"主推商品"必须出现在搜索前三。
  \item \textbf{新品冷启}：上架 7 天内的新品自然加权，让新品有曝光机会。
  \item \textbf{库存导向}：剩余 $\geq$ 100 件微加权，避免推无货 SKU。
  \item \textbf{客诉修正}：客服反馈"这个 query 总出错商品" $\to$ manual override。
\end{itemize}
\vspace{2mm}
路线图：\texttt{POST /admin/search/boost \{keyword, productID, weight\}}，
权重落 Redis，融合公式扩展为 \texttt{$\alpha \cdot kw + \beta \cdot sem + \gamma \cdot boost$}。
\srcnote{routes/router.go:130 admin/search/backfill 是同类型 admin 入口模板}
\end{frame}

\begin{frame}{客服话术：搜不到时怎么说}
\begin{tabular}{@{}lll@{}}
\toprule
业务码 & 含义 & 客服话术 \\
\midrule
\texttt{SE-EMBED-FAIL} & embedding 调用失败 & 智能搜索暂不可用，已切回关键词 \\
\texttt{SE-MILVUS-DOWN} & 向量库不可用     & （内部告警，对外无感）        \\
\texttt{SE-ES-DOWN}     & ES 不可用         & 搜索能力恢复中，请稍候           \\
\texttt{SE-EMPTY}       & 无结果           & 未找到，试试其他关键词           \\
\bottomrule
\end{tabular}
\vspace{3mm}
\begin{itemize}
  \item 业务码不暴露给前端，只在客服后台用。
  \item 用户侧统一文案"未找到"，避免技术名词刺激退出。
  \item gomall 的原则：\textbf{搜不到比搜错更可恢复}，宁可结果偏宽。
\end{itemize}
\srcnote{service/product.go:294--323 ES 失败 fall back DB 的降级模板}
\end{frame}

\begin{frame}{各角色关心的指标}
\begin{tabular}{@{}lll@{}}
\toprule
角色 & 关心 & 看哪个数字 \\
\midrule
产品 & GMV 转化     & 各 query 类型的搜索 $\to$ 详情 $\to$ 下单 \\
运营 & 主推命中率   & boost 商品在 top3 的出现率 \\
客服 & "搜不到" 客诉 & SE-* 业务码统计 \\
SRE  & 接口稳定性   & embedding / Milvus p95 + 错误率 \\
算法 & 召回 / 准确  & top-k recall / precision 离线评估 \\
\bottomrule
\end{tabular}
\vspace{2mm}
\small 算法的 recall / precision 是离线指标。\\
产品 / 运营关心的"实际收益"要看\textbf{在线 GMV 增量}，靠 AB 切桶量出来。
\end{frame}

\begin{frame}{回到业务承诺}
\begin{itemize}
  \item ES 单路：把"搜得到关键词"做到位。
  \item 加 Milvus：把"搜得到长尾意图"补上。
  \item 业务调权层：把"搜得到运营要卖的"留给非工程。
  \item 静默降级：把"搜索从不 5xx"的承诺延伸到三层。
\end{itemize}
\vspace{2mm}
\keypoint{从 LIKE 到 ES 是\textbf{召回率的一次跃迁}；
从 ES 到 ES + Milvus 是\textbf{召回率的第二次跃迁}。
两次跃迁解决的都不是"算法精度"，是\textbf{搜得到 vs 搜不到}的二元差距。}
\end{frame}

% ============================================================
\appendix
% ============================================================

\begin{frame}{业务侧 Q\&A（上）}
\textbf{Q1}：embedding 模型升级时，已有商品向量怎么迁？\\
\hspace{4mm}\small A：向量空间不可比，必须重算全量。做法是版本化 collection
（\texttt{product\_vector\_v1} / \texttt{v2}），后台跑全量 backfill 灌 v2，
搜索流量灰度切到 v2，验证后下掉 v1。整个过程对前端透明。

\vspace{2mm}
\textbf{Q2}：50/50 加权是不是可以做成可学习的？\\
\hspace{4mm}\small A：可以。两条路：离线 NDCG / MRR 网格搜，或在线 contextual bandit
按 query 类型动态调权。当前路线先把手动 50/50 上线，埋点 7 天再决定走哪条。

\vspace{2mm}
\textbf{Q3}：Milvus 挂了用户怎么感知？\\
\hspace{4mm}\small A：无感。客户端检测 \texttt{nopMilvusSearcher} 或网络挂 $\to$ 融合时
向量分数取 0 $\to$ 等价于纯 ES 单路。长尾词召回率掉，但主流量词不变。
告警 5 分钟内拉群，运营无须前台 banner。
\end{frame}

\begin{frame}{业务侧 Q\&A（下）}
\textbf{Q4}：以图搜图（跨模态）值不值得做？\\
\hspace{4mm}\small A：值得，但优先级低于文本语义。CLIP 类模型把图和文本 embed 到同一个空间，
用户拍照搜同款的转化率显著高于关键词。代价：图片 embed 慢，需要 GPU 推理服务。
路线图上排在文本语义之后。

\vspace{2mm}
\textbf{Q5}：要不要给每个用户建一个 user embedding 做个性化？\\
\hspace{4mm}\small A：典型推荐系统的路线。搜索侧用 user embedding 收益有限，因为搜索 query 已经表达了
即时意图。更划算的做法是把 user embedding 留给\textbf{推荐 / feed 流}，搜索专心做意图理解。

\vspace{2mm}
\textbf{Q6}：海外站点多语言搜索怎么扩？\\
\hspace{4mm}\small A：换 embedding 模型即可。multilingual-e5-large、bge-m3 这类多语种模型把
中 / 英 / 日 embed 到同一空间，无需为每种语言部署独立索引。ES 侧分词器要按语言切，
但向量侧一份 collection 就够。
\end{frame}

\begin{frame}{代码位置一览}
\textbf{main 已有（关键词层）}
\begin{description}
  \item[搜索路由入口]\srcnote{routes/router.go:40}
  \item[全量回填 admin 入口]\srcnote{routes/router.go:130}
  \item[搜索服务编排]\srcnote{service/search/service.go:11--16}
  \item[indexer 消费者]\srcnote{service/search/indexer.go:17--60}
  \item[全量回填实现]\srcnote{service/search/backfill.go:13--43}
  \item[ES 多字段查询 + 加权]\srcnote{repository/es/product\_index.go:156--209}
  \item[ProductSearch 降级]\srcnote{service/product.go:294--323}
  \item[ES 启动静默降级]\srcnote{cmd/main.go:60--69}
\end{description}

\vspace{2mm}
\textbf{PR 待合（语义层 + 融合）}
\begin{description}
  \item[Milvus 客户端]\srcnote{feat/milvus-vector：repository/milvus/milvus.go + product\_vector.go}
  \item[embedding 抽象]\srcnote{feat/semantic-search：service/search/embedding.go}
  \item[语义编排 + 融合]\srcnote{feat/semantic-search：service/search/semantic.go}
  \item[Milvus 接口桩]\srcnote{feat/semantic-search：service/search/milvus\_stub.go}
\end{description}
\end{frame}

\end{document}
